\section{Споріднені роботи}
\label{sec:related}

Social перебуває на перетині кризової інформатики, баз даних новинних подій, автоматизованої перевірки тверджень, генерації з опорою на пошук і практики цифрових розслідувань. Ці напрями працюють із різними одиницями: повідомленнями, зафіксованими подіями, твердженнями, відповідями або збереженими медіа. Тому їхні результати потребують різної валідації. Знайдена стаття може бути релевантною, але не підтримувати твердження; кілька повідомлень можуть мати одне походження; а декодоване відео може лишатися неперевіреним щодо місця, часу чи тлумачення. Це розрізнення використано для позиціювання дослідженого програмного забезпечення, а не для трактування суміжних систем як взаємозамінних конкурентів. У таблиці~\ref{tab:related} підсумовано їхні задокументовані призначення та межі порівняння.

Безпосередню постановку проблеми дає кризова інформатика. Імран та співавтори~\cite{imran2015crisis} організовують обробку екстрених повідомлень навколо фільтрації, класифікації, ранжування, агрегації, виділення та реферування. Така рамка пов'язує неоднорідний збір із потребою аналітика зменшити інформаційне перевантаження. Social використовує подібний погляд на робочий процес, але досліджує обмежену реконструкцію інциденту та придатні для перегляду артефакти. Його панелі платформ і правила добору для конкретного кейсу не становлять репрезентативної вибірки генеральної сукупності. Тому успішне виконання фільтра чи збирача не може встановити повноту, ситуаційну точність або корисність для служб реагування без окремої оцінки цих результатів.

Інфраструктури новинних подій показують, що впорядкування спостережень у події є усталеним завданням. GDELT відокремлює закодовані за CAMEO записи подій від згадок у новинах і надає впевненість виділення та географічні атрибути; його GKG містить сутності, теми й контекстні метадані на рівні документа~\cite{gdelt2015events,gdelt2015gkg}. ACLED спирається на кодування дослідниками та визначену таксономію політичних заворушень із явною точністю дати й місця та інформацією про джерела~\cite{acledCodebook}. Event Registry групує багатомовні новинні статті, виділяє атрибути подій і підтримує їх пошук та дослідження~\cite{leban2014eventregistry}. Ці роботи обґрунтовують явні одиниці подій і розрізнення повідомлень та самих подій. Social не претендує на новизну у виділенні чи кластеризації. Його вужче дослідницьке питання стосується того, як квитанції збору, доказові уривки, оцінки тверджень і гіпотези подій можуть залишатися придатно для перегляду пов'язаними в межах обмеженого розслідування.

Автоматизована перевірка фактів додає доповнювальну вимогу: оцінювати докази так само, як і передбачений висновок. FEVER поєднує твердження, виведені з Wikipedia, із добором доказів і мітками «підтримано», «спростовано» або «недостатньо інформації»~\cite{thorne2018fever}. AVeriTeC працює з твердженнями реального світу, використовуючи вебдокази, декомпозиції запитань і відповідей та обґрунтування, з окремою увагою до часових витоків і достатності доказів~\cite{schlichtkrull2023averitec}. InFact реалізує поетапний конвеєр перевірки фактів на основі пошуку та LLM~\cite{rothermel2024infact}. Це релевантні постановки завдань і прецеденти оцінювання. Контрактні перевірки та регресії Social для інциденту не дорівнюють результату на бенчмарках FEVER чи AVeriTeC, і жодного рейтингу точності щодо цих систем не передбачено. Правильність на рівні тверджень і якість реконструкції лишаються окремими від успіху програмної інтеграції.

Генерація з доповненням пошуком поєднує навчену генерацію із зовнішньою знайденою пам'яттю~\cite{lewis2020rag}. ALCE робить якість цитування окремим виміром оцінювання поряд із правильністю та плавністю~\cite{gao2023alce}. Це розрізнення важливе для згенерованих звітів розслідування: синтаксично наявна цитата не обов'язково релевантна, достатня чи вірна пов'язаному реченню. Посилання Social на джерела й уривки надають об'єкти, які можна перевірити, але саме їхнє існування не гарантує семантичної підтримки. Ані релевантність пошуку, ані словесну впевненість моделі не слід тлумачити як незалежно калібровану ймовірність істинності твердження про інцидент.

Обсерваторії поширення й інструменти перевірки медіа охоплюють інші частини процесу. Hoaxy відстежує обіг URL із моніторованих джерел і відповідних матеріалів перевірки фактів; його опис 2016 року прямо розмежовує відстеження й саму перевірку фактів~\cite{shao2016hoaxy}. InVID і WeVerify поєднують виділення ключових кадрів, зворотний пошук зображень, перегляд метаданих, криміналістичну допомогу та керовані інструменти аналітика~\cite{invidVerificationPlugin,weverifyVerificationPlugin}. Вони становлять вагомий попередній доробок у сфері інтегрованої допомоги з перевірки. Їхні задокументовані функції не перетворюють кожен діагностичний результат на фактичний вердикт. Так само локальне індексування медіа й кандидатні транскрипти Social слід відрізняти від незалежної автентифікації та перевірки контексту джерела.

Берклійський протокол розділяє збір, збереження, перевірку, слідчий аналіз і звітування в межах професійного й етичного процесу розслідування~\cite{berkeleyProtocol2022}. Це методичне керівництво, а не сертифікація програмного забезпечення. Практики дослідницького програмування щодо організованих даних, задокументованих перетворень, контролю версій і відтворюваних процесів дають доповнювальну інженерну рамку~\cite{wilson2017practices}. Внесок Social оцінено за конкретними реалізованими артефактами й обмеженими перевірками в межах цих усталених традицій. Ані використання термінології, ані проходження локальних тестів не встановлюють відповідності протоколу, допустимості як доказу, універсального охоплення платформ чи семантичної обґрунтованості для різних інцидентів. Такі ширші твердження потребують додаткових доказів і відповідного дизайну оцінювання.

\begin{table}[!htbp]
\centering
\footnotesize
\setlength{\tabcolsep}{3pt}
\renewcommand{\arraystretch}{1.12}
\caption{Якісне позиціювання за цитованими першоджерелами та цим аудитом коду. Клітинки описують задокументовані функції або обсяг валідації, а не виміряну перевагу. Давні статті та сторінки інструментів не встановлюють поточного доступу до платформ.}
\label{tab:related}
\begin{tabularx}{\textwidth}{@{}p{20mm}YYYYY@{}}
\toprule
Система & Основне призначення & Подання походження & Міркування про події & Невизначеність & Оцінювання / порівняння \\
\midrule
GDELT & Метадані новинних подій і контексту & Ідентифікатори подій і згадок; посилання на документи & Дії CAMEO; контекстні поля & Впевненість виділення; географічні оцінки & Кодова книга описує поля; спільного бенчмарка інциденту тут немає \\
ACLED & Закодовані події політичних заворушень & Джерела й примітки до подій & Таксономія подій, закодована дослідниками & Точність дати й місця; правила кодування & Методологія й процес перевірки; інша одиниця спостереження \\
Event Registry & Упорядкування новин у події & Пов'язані новинні статті & Багатомовна кластеризація й атрибути & Обсяг мов і кластеризації & Опис системи 2014 року; спільного тесту тут немає \\
Hoaxy & Обсерваторія поширення URL & Моніторовані URL і спостереження поширень & Аналіз обігу & Обрані джерела й вибірка платформи & Попередній аналіз поширення; без встановлення істинності \\
InVID--WeVerify & Допомога аналітику в перевірці медіа & Контекст, метадані й результати пошуку & Дослідження медіа під керівництвом аналітика & Діагностичні інструменти потребують тлумачення & Офіційні описи інструментів; спільного тесту продуктивності немає \\
Social & Обмежений процес роботи з доказами інциденту & Хеші збору, шляхи, уривки й типізовані зв'язки & Оцінки тверджень; гіпотези подій & Стани доступу й нерозв'язані докази & Локальні регресійні перевірки й перевірки цілісності; обмежена семантична валідація \\
\bottomrule
\end{tabularx}
\end{table}
